\begin{tikzpicture}[
  fig,
  font=\sffamily\scriptsize,
  every node/.style={align=center},
  lbl/.style={draw, thick, fill=white, inner sep=2.5pt},
  smalllbl/.style={draw, thick, fill=white, inner sep=2pt},
  wire/.style={thick},
  beam/.style={very thick},
  arrow/.style={->, thick},
  hit/.style={fill=black, draw=none},
]

% ============================================================
% GEOMETRY : 90 degree magnetic sector.
% The flight tube bends about the sector centre C = (3,-\Rb); an ion of the
% selected m/z travels on the nominal radius \Rb and reaches the detector,
% while heavier (larger radius) and lighter (smaller radius) ions are lost
% to the outer and inner wall respectively.
% ============================================================
\def\h{0.35}          % tube half-width
\def\Rb{2.2}          % nominal (selected m/z) radius
\def\Le{1.15}         % straight exit section after the bend
\pgfmathsetmacro{\Rout}{\Rb+\h}       % outer wall radius   (2.55)
\pgfmathsetmacro{\Rin}{\Rb-\h}        % inner wall radius   (1.85)
\pgfmathsetmacro{\Rpo}{\Rb+\h+0.35}   % outer pole radius   (2.90)
\pgfmathsetmacro{\Rpi}{\Rb-\h-0.35}   % inner pole radius   (1.50)
\pgfmathsetmacro{\Xex}{3+\Rb}         % exit-tube centreline x
\pgfmathsetmacro{\Yex}{-\Rb-\Le}      % exit-tube mouth y
\coordinate (C) at (3,-\Rb);          % sector centre

% ---------- ion source / acceleration block ----------
\draw[arrow] (-2.1,0.0) -- (-0.35,0.0);
\node[anchor=east] at (-2.15,0.0) {vaporised\\sample};
\draw[wire] (-0.35,-0.35) rectangle (2.15,0.35);
\draw[wire] (0.15,-0.05) .. controls (0.35,0.25) and (0.55,0.25) .. (0.75,-0.05);
\draw[wire] (0.75,-0.05) .. controls (0.55,-0.35) and (0.35,-0.35) .. (0.15,-0.05);
\foreach \x in {1.10,1.35,1.60}{\draw[wire] (\x,-0.35) -- (\x,0.35);}

% ============================================================
% ELECTROMAGNET : curved pole pieces flanking the bend
% ============================================================
\draw[Tan, line width=9pt, line cap=butt]
  ([shift={(78:\Rpo cm)}]C) arc[start angle=78, end angle=12, radius=\Rpo cm];
\draw[Tan, line width=9pt, line cap=butt]
  ([shift={(78:\Rpi cm)}]C) arc[start angle=78, end angle=12, radius=\Rpi cm];

% ============================================================
% FLIGHT TUBE : straight in, 90 degree bend, straight out
% ============================================================
\draw[wire] (2.15, \h) -- (3, \h)
  arc[start angle=90, end angle=0, radius=\Rout cm] -- (\Xex+\h, \Yex);
\draw[wire] (2.15,-\h) -- (3,-\h)
  arc[start angle=90, end angle=0, radius=\Rin cm] -- (\Xex-\h, \Yex);

% ============================================================
% ION STREAMS : selected m/z stays on the tube axis; the others do not
% ============================================================
% selected m/z -> follows the bend to the detector
\draw[beam, green!55!black] (2.15,0) -- (3,0)
  arc[start angle=90, end angle=0, radius=\Rb cm] -- (\Xex, \Yex);
% heavier ion: radius too large, under-deflected -> outer wall
\draw[beam, blue] (3,0)
  arc[start angle=90, end angle=38.63, radius=2.95cm] coordinate (hitH);
% lighter ion: radius too small, over-deflected -> inner wall
\draw[beam, red] (3,0)
  arc[start angle=90, end angle=8.14, radius=1.72cm] coordinate (hitL);
\fill[hit] (hitH) circle (1.6pt);
\fill[hit] (hitL) circle (1.6pt);

% ============================================================
% DETECTOR + ELECTRONICS
% ============================================================
\draw[wire] (\Xex-0.48,\Yex) rectangle (\Xex+0.48,\Yex-0.37);
\node[smalllbl, minimum width=1.5cm, minimum height=0.55cm]
  (amp) at (\Xex+2.35,\Yex-0.18) {amplifier\\ \& other electronics};
\draw[arrow] (\Xex+0.50,\Yex-0.18) -- (amp.west);

% ============================================================
% LABELS
% ============================================================
\node[lbl] (ionlab) at (0.2,2.05) {IONISATION};
\draw[arrow] (ionlab.south) -- (0.45,0.38);
\node[lbl] (acclab) at (2.9,2.05) {ACCELERATION};
\draw[arrow] (acclab.south west) -- (1.70,0.40);
\node[lbl] (magn) at (6.75,1.45) {ELECTROMAGNET};
\draw[arrow] (magn.south west) -- ([shift={(55:\Rpo cm)}]C);
\node[lbl] (deflab) at (7.45,-1.10) {DEFLECTION};
\draw[arrow] (deflab.west) -- ($(hitH)+(0.12,0)$);
\node[lbl] (detlab) at (3.45,\Yex-0.20) {DETECTION};
\draw[arrow] (detlab.east) -- (\Xex-0.50,\Yex-0.18);

\end{tikzpicture}
